% !PRINTED 500
\weberpage{500}
%<<CONT
11. \gesperrt{that, when the case of imprimitivity occurs, the group $P$ contains, in case 9. the whole symmetric group, in case 10. the whole alternating group, for each single system of imprimitivity.}

Repeated composition of a permutation with itself is denoted by
powers. If $\pi$ is any permutation, then among the successive
powers
\[
\pi^0, \pi^1, \pi^2, \pi^3 \ldots,
\]
in which $\pi^0 = 1$ denotes the identical permutation, the same
permutation must necessarily recur. But from $\pi^u = \pi^{u+e}$ it follows
that $\pi^e = 1$. If $e$ is the least positive exponent for which $\pi^e = 1$,
then the permutations
\weq{4}{1, \pi, \pi^2 \ldots \pi^{e-1}}
are all different from one another, while, as the sequence is continued with
$\pi^e, \pi^{e+1}, \ldots \pi^{2e-1} \ldots$, they repeat periodically in the same
order. The sequence (4) is called the \gesperrt{period of the permutation}
$\pi$, and $e$ is called the \gesperrt{degree of} $\pi$. The period
of $\pi$ forms a group, since composing any two of its elements
yields another element of the same period.
If $\pi$ is a cyclic permutation
\[
\pi = (0, 1, 2 \ldots n - 1),
\]
then for even $n$
\[
\pi^2 = (0, 2, 4 \ldots n - 2) (1, 3 \ldots n - 1),
\]
and for odd $n$
\[
\pi^2 = (0, 2, 4 \ldots n - 1, 1, 3 \ldots n - 2),
\]
and the degree of $\pi$ is in both cases equal to $n$.
If $\pi$ is decomposed into several cycles, the degree of $\pi$
is equal to the least common multiple of the number of terms of the single cycles, e.\,g.
\[
\pi = (0, 1) (2, 3, 4), \; \pi^2 = (0) (1) (2, 4, 3), \; \pi^3 = (0, 1) (2) (3) (4),
\]
% [sic] The print separates the disjoint cycles in $\pi^4$ and $\pi^5$ by commas, unlike $\pi^2$ and $\pi^3$ immediately above; the printed notation is preserved.
\[
\pi^4 = (0), (1), (2, 3, 4), \; \pi^5 = (0, 1), (2, 4, 3), \; \pi^6 = 1.
\]
If $\pi$ occurs in any permutation group $P$,
then $P$ contains the whole period of $\pi$.
